\documentclass[12pt,a4paper]{ctexart}
\usepackage{geometry}
\geometry{left=2.5cm,right=2.5cm,top=2.5cm,bottom=2.5cm}
\usepackage{amsmath}
\usepackage{hyperref}
\usepackage{enumitem}
\setlist{nosep}

\title{AI 工具使用详情}
\author{}
\date{2026年7月28日}

\begin{document}
\maketitle

% ==================== 领域知识与数据准备 ====================
\section*{领域知识与数据准备}

\subsection*{[AI-0-2]}

\textbf{所用AI工具名称和版本：}\par
DeepSeek, DeepSeek-V4, 深度求索（DeepSeek）

\textbf{具体使用目的和环节：}\par
赛题分析与领域知识检索。团队在接到赛题后，需要系统了解评分者信度（ICC, Kendall's W, Fleiss' Kappa）、评委偏差分类框架（严厉度/宽容度、中心化趋势、光环效应）、相关性vs一致性的概念区分等教育测量学领域知识，以及通用的数学建模方法（熵权法、TOPSIS、K-Means、Kruskal-Wallis、Isolation Forest等），以指导后续建模方案的设计。团队确定上述知识领域后，要求AI通过联网检索整理领域知识卡片，并在读题阶段进行头脑风暴，梳理子问题间的依赖关系，给出初步方法论建议。

\textbf{关键交互记录（重要提示词与回复）：}\par

\begin{quote}
提示词：这是一道关于竞赛网评结果建模与分析的C题。请先检索以下领域知识：评分者信度（ICC, Kendall W, Fleiss Kappa）、评委偏差分类框架（严厉度/宽容度、中心化趋势、光环效应）、多组比较方法（Kruskal-Wallis），整理成知识卡片，包含阈值标准和适用场景。然后阅读赛题，梳理5个子问题之间的依赖关系，给出初步方法论建议。
\end{quote}

\textit{AI回复：} 整理了ICC模型选择指南（Koo \& Li, 2016）、评分者信度三维框架（Stemler, 2004）、评委偏差分类框架、多组比较方法等知识卡片。识别出子问题依赖关系：Q2→Q3→Q4为一条线，Q1→Q5为另一条线。提出了Spearman+Pearson两步走、四维度指标体系（信度/效度/公平性/区分力）、熵权法+TOPSIS综合评价、Kruskal-Wallis+Dunn检验等方法论建议。核心发现：评委完全不跨题，五题使用独立评委池，因此Q4中题目效应和评委池效应无法分离，只能报告"存在差异"。

\textbf{采纳和人工修改情况：}\par
团队审核了全部知识卡片的阈值标准和文献来源，采纳了ICC(2,1)成对比较、Kruskal-Wallis+Dunn检验等方法建议。子问题依赖关系图被直接用于后续工作安排。方案设计中团队补充了关键约束——Q4中题目效应与评委池效应无法分离，并要求在论文中坦诚此局限性。部分方法（如Fleiss' Kappa）因数据为连续评分而非分类数据，经团队评估后不予采用。

\subsection*{[AI-0-3]}

\textbf{所用AI工具名称和版本：}\par
DeepSeek, DeepSeek-V4, 深度求索（DeepSeek）

\textbf{具体使用目的和环节：}\par
数据清洗与数据画像。团队在建模前需要对附件A-E五个题目的评审数据进行系统性清洗和探索性分析，包括：检测重复记录与缺失值、统计各题数据量及入围比例、识别异常评分模式；随后进行多路径数据画像（聚类、降维、异常检测），为建模方案提供数据结构层面的验证。

\textbf{关键交互记录（重要提示词与回复）：}\par

\begin{quote}
提示词：读取附件A-E.xlsx，合并为统一数据集。检查重复记录（按阅卷号）、缺失值分布、各题数据量和入围比例。对评分列做IQR异常值检测，但异常值不做删除。统计评委跨题情况。输出清洗报告。

清洗完成后，对196位评委的11维特征矩阵做三条路径的数据画像：路径A——K-Means++聚类（K由肘部法则确定），路径B——PCA降维（前3个主成分），路径C——Isolation Forest异常检测。三条路径交叉验证，输出画像报告。
\end{quote}

\textit{AI回复：} 清洗结果：4876行×12列，零重复记录，成绩列2802条NaN为正常淘汰。IQR异常评分占比0.1\%-2.3\%，不做删除。确认196位评委完全不跨题。聚类路径发现两个评委群体（"严格型"132人，"宽松型"64人），Silhouette=0.283。PCA路径前3主成分累积解释85.8\%方差：PC1为评分水平因子（54.7\%）、PC2为评分离散度因子（22.0\%）、PC3为获奖率因子（9.1\%），PC3与PC1正交说明标准分公式有效消除了尺度差异。异常检测路径发现27位异常评委（13.8\%），分为极端严格（z<-1.5）、极端宽松+低离散（z>+2.0且$\sigma$低）、高离散度三类。三条路径交叉验证一致。

\textbf{采纳和人工修改情况：}\par
团队审核了清洗报告的数据准确性，确认异常评分保留不删除的策略——这些是评委评分差异的正常体现而非数据错误。数据画像结果直接支撑了后续方案设计：PCA揭示的"获奖率独立于评分水平"验证了标准分公式的有效性；聚类揭示的"严格/宽松"两极分化强化了Q2中"公平性"维度的必要性；异常检测揭示的三类异常评委为Q3/Q5的异常评委标记和降权策略提供了依据。

% ==================== 子问题 1 ====================
\section*{子问题 1：网评与终评相关性分析}

\subsection*{[AI-1-1]}

\textbf{所用AI工具名称和版本：}\par
DeepSeek, DeepSeek-V4, 深度求索（DeepSeek）

\textbf{具体使用目的和环节：}\par
子问题1方案设计与代码实现。团队确定了"Spearman+Pearson两步走+多维度评估"的总体框架后，要求AI完成以下工作：(1) 实现按评委ID分组计算z-score（而非按列位置），取4位评委z-score均值作为网评标准分；(2) 分五题计算Spearman $\rho$和Pearson r，报告效应量和置信区间；(3) 构建前55\%筛选的混淆矩阵，计算命中率和假阴性率；(4) 绘制ROC曲线，计算AUC评估网评对一等奖的区分能力。

\textbf{关键交互记录（重要提示词与回复）：}\par

\begin{quote}
提示词：对Q1实现标准分按评委计算——建立judge\_mu\_sigma映射，每个评委用自己的$\mu$和$\sigma$做z-score，取4个z-score均值。然后分五题计算Spearman $\rho$（含95\% CI和效应量），对入围论文计算Pearson r。构建55\%筛选混淆矩阵，计算命中率和假阴性率。绘制全题汇总ROC曲线。代码文件命名为c2\_q1q4\_model.py。
\end{quote}

\textit{AI回复：} 代码实现了按评委ID的z-score计算（c2\_q1q4\_model.py, L34-50），通过judge\_mu\_sigma字典建立评委到均值/标准差的映射。Spearman $\rho$全题汇总为0.797（强相关），AUC=0.932（卓越），前55\%筛选命中率74.8\%，假阴性率3.3\%。C题表现最弱：$\rho$=0.762、r=0.319、AUC=0.872（仍属良好）。

\textbf{采纳和人工修改情况：}\par
团队审查代码后发现初始版本存在关键错误——按列位置而非按评委ID计算z-score，导致不同评委的评分被混合计算均值和标准差。团队要求AI修复此问题后重新运行全部Q1计算。修复后Spearman $\rho$从0.737升至0.797，AUC从0.893升至0.932，网评有效性结论更加稳固。团队进一步要求AI对C题三指标联动（$\rho$高而r低+低AUC）给出实际含义解读，确认了"C题网评初筛功能正常但入围论文内部分数拉不开"的结论。

\subsection*{[AI-1-2]}

\textbf{所用AI工具名称和版本：}\par
硅基流动图像理解模型, Qwen3-VL-30B, 阿里云 / 硅基流动（SiliconFlow）

\textbf{具体使用目的和环节：}\par
Q1相关图表的可视化审查。子问题1生成了q1-spearman-bar、q1-confusion-heatmap、q1-roc五题、q1-top55-bar等图表。团队要求AI利用多模态模型逐一检查图表的标注正确性、视觉可读性和数据标注准确性，并通过自动修复流程迭代至通过标准。审查内容包括：横纵轴标签是否完整、图例是否遮挡数据、标题是否准确、数值标注是否与计算输出一致。

\textbf{关键交互记录（重要提示词与回复）：}\par

\begin{quote}
提示词：请依次审查以下图表的质量、标注正确性、数据解读和建模适用性——q1-spearman-bar.png, q1-confusion-heatmap.png, q1-roc.png, q1-top55-bar.png。检查横纵轴标签、图例位置、数值标注。
\end{quote}

\textit{AI回复：} q1-spearman-bar通过（轴标签"Spearman $\rho$"正确，数值标注与模型输出一致）。q1-confusion-heatmap通过（四个象限标注清晰，颜色映射合理）。q1-roc通过（五条ROC曲线可辨识，AUC标注完整）。q1-top55-bar通过（各题命中率/假阴性率柱状图清晰，高于55\%基准线标注完整）。4/4图表通过多模态审查。

\textbf{采纳和人工修改情况：}\par
团队审核了多模态审查结论，确认所有Q1图表标注正确、数据一致。特别核实了q1-roc图中五条ROC曲线的AUC数值与终端输出一致，q1-confusion-heatmap中混淆矩阵的四个象限数值与原始计算结果匹配。

% ==================== 子问题 2 ====================
\section*{子问题 2：评委基本素质指标体系}

\subsection*{[AI-2-1]}

\textbf{所用AI工具名称和版本：}\par
DeepSeek, DeepSeek-V4, 深度求索（DeepSeek）

\textbf{具体使用目的和环节：}\par
子问题2指标体系设计与指标计算实现。团队参考Stemler(2004)的三维评分者信度框架，结合竞赛实际场景，将指标体系扩展为四维度（信度/效度/公平性/区分力）。团队确定各维度的量化指标后，要求AI实现：(1) 成对ICC(2,1)计算每位评委的信度；(2) 评分与最终奖项的Spearman $\rho$计算效度；(3) 评分均值偏离度的z值计算公平性；(4) 评分标准差计算区分力。同时要求AI验证四维度间的独立性（交叉相关性），确保指标体系无严重冗余。

\textbf{关键交互记录（重要提示词与回复）：}\par

\begin{quote}
提示词：对196位评委构建四维度指标体系。信度：成对ICC(2,1)均值，使用pingouin或手写ANOVA分解实现。效度：评委评分与奖项等级的Spearman $\rho$，需包含淘汰论文（奖项填0）。公平性：评委均值偏离同题全体均值的z值绝对值。区分力：评委自身评分的标准差。每维度分题目内部归一化到[0,1]。验证四维度交叉相关性，确认无严重冗余。
\end{quote}

\textit{AI回复：} 实现了pairwise\_icc()函数（c2\_q2q3\_model.py），通过ANOVA分解计算MS\_between和MS\_error，公式为(MS\_between-MS\_error)/(MS\_between+MS\_error)。效度计算中淘汰论文奖项填充为0。四维度交叉验证：最大Spearman r=0.32（效度×区分力），公平性与其他维度完全独立（r<0.12），评分均值与获奖率r=0.10（标准分有效）。无严重冗余，指标体系结构合理。

\textbf{采纳和人工修改情况：}\par
团队审查后发现初始实现中两个关键问题：(1) 信度维度误用Spearman相关系数替代ICC，经团队指出后AI改为正确的ICC(2,1)实现——Spearman只衡量秩一致性但忽略绝对评分偏移，ICC同时捕捉趋势一致和绝对值一致；(2) 效度计算仅包含入围论文（约42\%），团队要求纳入淘汰论文（奖项填充为0），使得效度能完整衡量评委"区分入围vs淘汰"的能力。两项修复后团队要求重跑Q2→Q3→Q4→Q5全部代码和图表。修复后ICC信度值系统性低于原Spearman值，更真实反映了评委间绝对评分一致性的难度。团队成员确认所有修复后的指标值在业务逻辑上合理。

% ==================== 子问题 3 ====================
\section*{子问题 3：评委素质数学模型}

\subsection*{[AI-3-1]}

\textbf{所用AI工具名称和版本：}\par
DeepSeek, DeepSeek-V4, 深度求索（DeepSeek）

\textbf{具体使用目的和环节：}\par
子问题3的熵权法+TOPSIS综合评价模型实现。团队决定采用客观赋权方法（熵权法）避免主观赋权引入的偏差，以TOPSIS进行多指标综合评价，并用K-means++对评委进行聚类分层。团队要求AI实现：(1) 对四维度指标正向化+Min-Max归一化；(2) 熵权法计算各维度权重；(3) TOPSIS计算每位评委的综合贴近度；(4) K-means++聚类（K由肘部法则+Silhouette系数确定，业务约束K$\geq$3），将评委分为优秀/合格/需关注三档；(5) 等权TOPSIS作为稳健性检验对照。

\textbf{关键交互记录（重要提示词与回复）：}\par

\begin{quote}
提示词：对五题分别建模。Q2输出的四维度指标矩阵做Min-Max归一化后，用熵权法赋权：Pij=zij/$\sum$zij, ej=-1/ln(K)·$\sum$Pij·ln(Pij), wj=(1-ej)/$\sum$(1-ej)。用TOPSIS计算贴近度Si=Di-/(Di++Di-)。对TOPSIS得分做K-means++聚类，K由肘部法则+Silhouette系数联合确定，业务约束K$\geq$3。输出五题熵权分布、聚类分层结果、Silhouette系数。用等权TOPSIS做稳健性检验。
\end{quote}

\textit{AI回复：} 代码实现了完整熵权+TOPSIS流程（c2\_q2q3\_model.py）。熵权分布：区分力维度权重最高（0.304-0.442），公平性维度权重最低（0.145-0.249）——区分力权重高是因评委间标准差变异最大，属数据真实反映。TOPSIS得分范围分布合理。聚类结果：K=3（部分题目K=2因评委数不足，已修正下限为K$\geq$3），Silhouette系数0.241-0.329，分层结构合理。稳健性检验：熵权TOPSIS与等权TOPSIS的Spearman r=0.851-0.971，方法稳健。输出图表包括q3-weight-bar、q3-topsis-dist、q3-cluster-scatter、q3-robust-scatter。

\textbf{采纳和人工修改情况：}\par
团队审查后确认：(1) 熵权法区分力权重最高的结果与PCA的分析一致——评委评分离散度的差异确实最大，这不算法缺陷而是数据结构特征；(2) K-means++聚类在E题曾出现K=2（Silhouette=0.329），团队要求AI在代码中添加K$\geq$3的业务约束后修复；(3) 团队要求AI为每个聚类结果给出实际含义的解读，如Silhouette=0.243意味着"评委在各维度的表现差异较大，分层边界不够清晰但不影响宏观分组意义"。

% ==================== 子问题 4 ====================
\section*{子问题 4：不同题目评委表现差异分析}

\subsection*{[AI-4-1]}

\textbf{所用AI工具名称和版本：}\par
DeepSeek, DeepSeek-V4, 深度求索（DeepSeek）

\textbf{具体使用目的和环节：}\par
子问题4的多组比较分析实现。团队确定使用非参数Kruskal-Wallis检验（因196位评委分组后各题样本量差异大——A题仅14位，不满足ANOVA正态假设），并补充Dunn事后检验定位差异来源。团队要求AI：(1) 以Q3输出的TOPSIS得分为因变量、题目为自变量做KW检验；(2) 报告$\eta^2$效应量；(3) 做Dunn事后检验（Bonferroni校正×10）；(4) 额外对信度、效度、区分力三个维度分别做Dunn检验，定位综合得分差异的来源维度。

\textbf{关键交互记录（重要提示词与回复）：}\par

\begin{quote}
提示词：用scipy.stats.kruskal对五题TOPSIS得分做Kruskal-Wallis检验，计算$\eta^2$效应量。做Dunn事后检验，z检验公式自行实现，p值Bonferroni校正×10。此外对信度、效度、区分力三个子维度分别做Dunn检验，定位题目间差异来源于哪个维度。代码放在c2\_q1q4\_model.py的Q4部分。
\end{quote}

\textit{AI回复：} KW检验结果：H=11.710, p=0.019, $\eta^2$=0.040（效应量中等）。Dunn事后检验发现B-D题对p\_adj=0.031显著，B-E题对p\_adj=0.030显著，其他8对不显著。分维度Dunn检验揭示：信度（H=53.1, p<0.001）和效度（H=41.0, p<0.001）在五题间差异极显著，区分力（H=20.3, p<0.001）显著，公平性（H=2.8, p=0.59）无显著差异。信度维度的差异主要来自B题与E题的对比，效度维度上A题和E题评委表现更优。公平性五题一致说明标准分公式在各题均有效。

\textbf{采纳和人工修改情况：}\par
团队审核了统计检验结果，重点关注了公平性五题无显著差异的发现——这验证了标准分公式在消除评委尺度差异方面五题表现一致。团队要求AI补充分维度Dunn检验的代码，以精确定位综合得分的差异来源。团队在论文中坦诚了核心局限：由于评委完全不跨题，题目效应与评委池效应无法分离，本文只能报告"存在差异"而非建立因果。

% ==================== 子问题 5 ====================
\section*{子问题 5：网评成绩利用策略}

\subsection*{[AI-5-1]}

\textbf{所用AI工具名称和版本：}\par
DeepSeek, DeepSeek-V4, 深度求索（DeepSeek）

\textbf{具体使用目的和环节：}\par
子问题5的敏感性分析与改进建议。团队综合Q1-Q4结论后，决定从两个路径评估网评的利用策略：(1) 模拟网评权重$\alpha$从0到1变化时，网评-终评混合排名与奖项的Spearman $\rho$变化（敏感性分析）；(2) 基于Q3的评委TOPSIS得分对评委做素质加权，对比等权网评与加权网评的Spearman $\rho$提升。团队要求AI实现这两条路径并基于结果提出改进建议。

\textbf{关键交互记录（重要提示词与回复）：}\par

\begin{quote}
提示词：实现Q5两条分析路径。路径一：敏感性分析——$\alpha$从0到1（步长0.1），计算$\alpha$·rank(z\_mean)+(1-$\alpha$)·rank(award)与奖项的Spearman $\rho$，绘制$\rho$-$\alpha$曲线。路径二：素质加权——w=TOPSIS得分/$\sum$TOPSIS得分，计算加权网评标准分，对比等权和加权的Spearman $\rho$。列出降权建议：TOPSIS分层为'需关注'的评委名单。代码命名为c2\_q5\_model.py。
\end{quote}

\textit{AI回复：} 敏感性分析曲线显示：五题$\rho$均随$\alpha$增大而单调下降，$\alpha$=0（纯终评）时$\rho$最高但终评不可观测，$\alpha$=0.25（当前赛制等权）时$\rho$介于0.64-0.81之间。C题$\rho$从$\alpha$=0到$\alpha$=0.25下降最陡（-0.12），印证Q1中C题网评质量最弱的结论。素质加权后：加权网评的$\rho$在所有五题上均高于等权网评（提升幅度0.006-0.218），C题提升最为显著（+0.218）。降权建议：共有23-46位评委的TOPSIS得分位于最低档，建议对这部分评委的网评评分降权或不纳入最终成绩。

\textbf{采纳和人工修改情况：}\par
团队审查后发现初始实现使用了"秩混合"替代"分值混合"——这说明$\alpha$变化对Spearman $\rho$的影响被人为压缩。团队与AI讨论后确认：由于集中评审的$\bar{z}_{\text{central}}$不可观测，秩混合是不得已的近似方案。团队要求AI在论文中明确标注此简化假设及其局限性。此外，团队要求AI区分两个不同的"等权"概念：Q5中的等权是当前赛制实际使用的评委简单平均（基线对照），Q3中稳健性检验的等权是等权TOPSIS（方法稳健性验证）——两者目的不同，避免读者混淆。

% ==================== 跨子问题审查 ====================
\section*{跨子问题代码审查与图表生成}

\subsection*{[AI-0-4]}

\textbf{所用AI工具名称和版本：}\par
DeepSeek, DeepSeek-V4, 深度求索（DeepSeek）

\textbf{具体使用目的和环节：}\par
跨子问题一致性审查。Q1-Q5五个子问题的代码实现完成后，团队要求AI执行两轮系统性审查：(1) 代码深度审计——检查各子问题代码与数学推导的一致性、代理值选择的正确性、方案确认书的匹配度；(2) 跨子问题审查——检查符号体系的统一性、子问题间数据流衔接的正确性、核心假设的一致性。审查发现的代码问题需修复后重跑全部模型和图表。

\textbf{关键交互记录（重要提示词与回复）：}\par

\begin{quote}
提示词：请审查所有代码文件（c2\_q1q4\_model.py, c2\_q2q3\_model.py, c2\_q5\_model.py, c3\_generate\_charts.py）与数学推导（math-sub1.tex到math-sub5.tex）的一致性。重点检查：(1) z-score是按评委还是按列计算？(2) 信度用Spearman还是ICC？(3) 效度是否纳入了淘汰论文？(4) Q5是分值混合还是秩混合？(5) K-means的K是否满足K$\geq$3约束？(6) 内部报告的数值是否已过期？
\end{quote}

\textit{AI回复：} 发现6个问题：(1) z-score按列计算而非按评委——混合了不同评委的评分（关键，已修复）；(2) 信度用Spearman替代ICC——高估了信度值（关键，已修复）；(3) 效度仅用入围论文——排除了淘汰论文的信息（关键，已修复）；(4) Q5用秩混合替代分值混合——低估了$\alpha$敏感度（已标注）；(5) E题K=2违反K$\geq$3约束（已修复）；(6) 内部报告数值全部过时（已刷新）。18项审查中有12项通过。全部修复后代码审查通过。

\textbf{采纳和人工修改情况：}\par
团队逐项验证了AI审查发现的每个问题，确认修复后的代码与数学推导完全一致。特别核实了：(1) z-score按评委ID计算后Spearman $\rho$从0.737升至0.797；(2) ICC替代Spearman后信度值更真实但系统性降低——这符合预期因为ICC同时捕捉绝对偏差；(3) 效度纳入淘汰论文后更多评委的得分被重新校准。修复后的全部图表由AI利用多模态模型逐张审查，20张图表全部通过。

\subsection*{[AI-0-5]}

\textbf{所用AI工具名称和版本：}\par
DeepSeek, DeepSeek-V4, 深度求索（DeepSeek）；硅基流动图像理解模型, Qwen3-VL-30B, 阿里云 / 硅基流动（SiliconFlow）

\textbf{具体使用目的和环节：}\par
论文撰写与图表生成。全部建模计算完成后，团队要求AI完成：(1) 批量生成全部20张图表（c3\_generate\_charts.py），图表标注中文化、图例位置调整避免遮挡、标题精简；(2) 利用多模态模型逐张审查图表质量（20/20通过）；(3) 生成各子问题的数学推导LaTeX文件（math-sub1.tex至math-sub5.tex）；(4) 撰写论文正文main.tex并编译为PDF；(5) 多轮精修——补充13处数据结果的实际竞赛场景含义解读（如C题三指标联动分析、熵权极端值的实际意义、Silhouette系数的分层含义等），使论文分析从技术性描述转向实际应用解读。

\textbf{关键交互记录（重要提示词与回复）：}\par

\begin{quote}
提示词：批量生成全部图表。横轴标签统一为"A题""B题"等中文格式，图例位置设为upper left避免遮挡，标题去除冗余括号注释。生成后调用多模态模型逐张审查，检查标注正确性、视觉可读性、数据与终端输出的一致性。编译论文main.tex，检查编译错误。

论文中有很多地方只给了知识性技术性说明但没有给出解读。比如：对四个维度分别做Kruskal-Wallis检验的H值和p值都报了，但没有说这些差异在实际竞赛场景中意味着什么。请系统性发现此类缺失，并予以解读——每条数据结果都要对应到竞赛组织的实际意义。
\end{quote}

\textit{AI回复：} 20张图表全部生成并逐张通过多模态审查（"通过"级别）。论文编译成功（24页，无错误）。发现并修复了13处解读缺失：(1) C题三指标联动——$\rho$高而r低意味着网评初筛功能正常但入围论文内部排序不可靠；(2) B题区分力权重最高（0.442）——B题评委给分习惯差异极大，标准分校正对B题必要且有效；(3) 公平性五题无差异——标准分公式消除尺度差异的效果五题一致；(4) E题区分力权重最低——不是评委表现差，而是E题评委打分尺度本来就比较统一，区分力维度区分不开；(5) 另9处涉及各个具体指标的解读。论文页数从22页增加至24页。

\textbf{采纳和人工修改情况：}\par
团队审核了全部20张图表的多模态审查结果，确认通过。论文解读修复方面，团队特别关注了"技术性vs实际性"的区分——要求每条解读必须映射到竞赛组织的实际操作场景。例如"B题区分力权重最高"的技术性解读是"B题评委间标准差变异最大"，而实际性解读是"这意味着标准分校正对B题的必要性远高于其他题：如果不做z-score，那些'敢打高分'的评委将主导B题排名"。团队逐一核实了13处解读的准确性和实际相关性后定稿。

% ==================== 固定条目 ====================
\section*{终稿编译}

\subsection*{[AI-0-0]}

\textbf{所用AI工具名称和版本：}\par
DeepSeek, DeepSeek-V4, 深度求索（DeepSeek）

\textbf{具体使用目的和环节：}\par
终稿 LaTeX 编译与排版。论文正文（含模型描述、公式推导、结果分析）由参赛队独立撰写，AI 仅将已完成的文字、公式和图表转化为 LaTeX 代码并编译为 PDF，全程不参与内容生成。

\textbf{关键交互记录（重要提示词与回复）：}\par

\begin{quote}
提示词：将以下论文正文编译为 LaTeX，不修改任何实质内容。
\end{quote}

\textit{AI 回复：}（编译完成，输出 PDF）

\textbf{采纳和人工修改情况：}\par
AI 编译后人工检查排版效果，修正格式问题，内容未做任何 AI 修改。

\subsection*{[AI-0-1]}

\textbf{所用AI工具名称和版本：}\par
DeepSeek, DeepSeek-V4, 深度求索（DeepSeek）

\textbf{具体使用目的和环节：}\par
本文件的条目整理、格式编排与 LaTeX 编译。团队导出竞赛期间的全部对话记录后，要求AI按照竞赛规定的条目格式汇总其中涉及的AI工具使用情况并编译为支撑材料 PDF。

\textbf{关键交互记录（重要提示词与回复）：}\par

\begin{quote}
提示词：按照竞赛规定的[AI-X-Y]条目格式，汇总四个子问题中所有AI工具使用情况，每个条目需包含工具名称版本、使用目的和环节、关键交互记录、采纳和人工修改情况，编译为PDF。
\end{quote}

\textit{AI 回复：}（完成全部条目的整理与编译）

\textbf{采纳和人工修改情况：}\par
团队审核了全部条目的内容完整性和表述准确性，确认后定稿。

\end{document}
